\documentclass[11pt]{article}

\usepackage[margin=1in]{geometry}
\usepackage{amsmath}
\usepackage{amssymb}
\usepackage{amsthm}
\usepackage{hyperref}
\usepackage{enumitem}
\usepackage{color}

\title{Comprehensive Problem Set: Separable Differential Equations}
\author{Math Curated Problem Sets}
\date{\today}

\begin{document}

\maketitle
\tableofcontents
\newpage

% ==========================================
% CONCEPT CHECKLIST & QUESTION MAPPING
% ==========================================
\section{Concept Checklist \& Question Mapping}
This section categorizes all required concepts and problem types extracted from the provided text, mapped directly to the problems created in this document. It serves as a study guide and a verification that all techniques have been rigorously tested.

\begin{itemize}
    \item \textbf{1. Identification and Setup}
    \begin{itemize}
        \item Determining if an ODE is separable (Algebraic manipulation, blocked by addition/subtraction): \textbf{Q1, Q2, Q3}
    \end{itemize}
    
    \item \textbf{2. General Implicit \& Explicit Solutions}
    \begin{itemize}
        \item Basic variable separation and integration: \textbf{Q4, Q9}
        \item Exponentiating Logarithms to Isolate $y$ (The absolute value trick $A = \pm e^C$): \textbf{Q8, Q10, Q13}
    \end{itemize}
    
    \item \textbf{3. Key Algebraic \& Calculus Manipulations}
    \begin{itemize}
        \item Splitting Exponentials ($e^{a+b} = e^a e^b$): \textbf{Q2, Q5}
        \item Factoring to uncouple variables: \textbf{Q3, Q6, Q8, Q10}
        \item Integration by Parts (LIATE rule): \textbf{Q5, Q11, Q15}
        \item Trigonometric Identities (e.g., $1/\cos^2(y) = \sec^2(y)$): \textbf{Q7, Q12}
        \item Polynomial Long Division in ODEs: \textbf{Q10}
    \end{itemize}
    
    \item \textbf{4. Initial Value Problems (IVPs)}
    \begin{itemize}
        \item Solving for $C$ immediately after integration: \textbf{Q11, Q12, Q13}
        \item Proper square root branch selection (Positive vs. Negative): \textbf{Q14, Q30}
    \end{itemize}
    
    \item \textbf{5. Common Variations \& Untested Concepts}
    \begin{itemize}
        \item Singular Solutions (Lost constant solutions from dividing by zero): \textbf{Q16, Q17, Q18}
        \item Partial Fractions Integration: \textbf{Q18, Q20, Q21}
        \item The Logistic Growth Equation: \textbf{Q19, Q25}
    \end{itemize}
    
    \item \textbf{6. Real-World Application and Modeling}
    \begin{itemize}
        \item Continuous Wealth Accumulation (Compound Interest): \textbf{Q22}
        \item Price Elasticity of Demand: \textbf{Q23}
        \item The Vasicek Interest Rate Model (Mean-reverting): \textbf{Q24}
    \end{itemize}
    
    \item \textbf{7. Advanced Diagnostic Testing (''Find the Flaw'')}
    \begin{itemize}
        \item Flaw: Constant of Integration Exponentiation: \textbf{Q26}
        \item Flaw: Misusing Natural Log on Power Rule Integrals: \textbf{Q27}
        \item Flaw: Fake Separation (Separating via addition/subtraction): \textbf{Q28}
        \item Flaw: Multiplying instead of Dividing to move variables: \textbf{Q29}
        \item Flaw: Incorrect Square Root Branching: \textbf{Q30}
    \end{itemize}
\end{itemize}

\newpage

% ==========================================
% PROBLEM SETS
% ==========================================
\section{Problem Sets}

\subsection*{Part 1: Identification and Setup}
\textbf{Problem 1:} Determine whether the given differential equation is separable. Show your reasoning.
\[ \frac{dy}{dx} + 4\ln(x-y) = 0 \]

\textbf{Problem 2:} Determine whether the given differential equation is separable. Show your reasoning.
\[ \frac{dy}{dx} = \frac{x^2 e^{3x-4y}}{x^4 + 1} \]

\textbf{Problem 3:} Determine whether the given differential equation is separable. Show your reasoning.
\[ (ty^3 + 5y^3)dy - 7t \, dt = 0 \]

\subsection*{Part 2: General Implicit \& Explicit Solutions}
\textbf{Problem 4:} Solve the differential equation implicitly.
\[ \frac{dx}{dt} = 5xt^4 \]

\textbf{Problem 5:} Solve the differential equation. (Implicit solution is acceptable).
\[ \frac{dy}{dt} = \frac{3}{y e^{t+3y}} \]

\textbf{Problem 6:} Solve the differential equation implicitly.
\[ (t^2 + 4t^2 y^3) \, dt + e^{t^3} y^2 \, dy = 0 \]

\textbf{Problem 7:} Solve the differential equation explicitly for $y$.
\[ \frac{dy}{dx} = 9x \cos^2(y) \]

\textbf{Problem 8:} Solve the differential equation explicitly for $y$.
\[ \frac{dy}{dx} = x^4(y - 2) \]

\textbf{Problem 9:} Solve the differential equation implicitly.
\[ \sqrt{y} \, dx + (3+x) \, dy = 0 \]

\textbf{Problem 10:} Solve the differential equation explicitly for $y$.
\[ \frac{dy}{dx} = \frac{x^2y - 2y}{x+5} \]

\subsection*{Part 3: Initial Value Problems (IVPs)}
\textbf{Problem 11:} Solve the initial value problem implicitly.
\[ \frac{1}{t} \frac{dy}{dt} = \frac{y \cos(t)}{y^3 + 1}, \quad y(\pi) = 1 \]

\textbf{Problem 12:} Solve the initial value problem explicitly for $y$.
\[ \frac{dy}{dx} = 2x \cos^2(y), \quad y(1) = \frac{3\pi}{4} \]

\textbf{Problem 13:} Solve the initial value problem explicitly for $y$.
\[ \frac{dy}{dx} = x^3(y + 5), \quad y(0) = 2 \]

\textbf{Problem 14:} Solve the initial value problem explicitly for $y$.
\[ \frac{dy}{dx} = -\frac{x}{y}, \quad y(4) = -3 \]

\textbf{Problem 15:} Solve the initial value problem implicitly.
\[ \frac{dy}{dt} = t \ln(t) y^2, \quad y(1) = -2 \]

\subsection*{Part 4: Singular Solutions}
\textbf{Problem 16:} Find the general explicit solution. Identify any singular solutions that are lost during variable separation.
\[ \frac{dy}{dx} = y^{1/3} \]

\textbf{Problem 17:} Find the general explicit solution. Identify any singular solutions.
\[ \frac{dy}{dx} = x(y-5)^{3/4} \]

\textbf{Problem 18:} Find the general implicit solution. Identify any singular solutions.
\[ \frac{dy}{dx} = 16 - y^2 \]

\subsection*{Part 5: Logistic Growth \& Partial Fractions}
\textbf{Problem 19:} Solve the logistic growth equation implicitly.
\[ \frac{dy}{dt} = 3y\left(1 - \frac{y}{10}\right) \]

\textbf{Problem 20:} Solve the differential equation implicitly.
\[ \frac{dy}{dx} = \frac{y^2 - 1}{x} \]

\textbf{Problem 21:} Solve the differential equation explicitly for $y$.
\[ \frac{dy}{dt} = y^2 - 7y + 12 \]

\subsection*{Part 6: Real-World Application and Modeling}
\textbf{Problem 22 (Continuous Wealth):} A portfolio's wealth $W(t)$ grows continuously at an interest rate $r = 0.06$. Formulate the differential equation. Given an initial investment of $\$5000$ at $t=0$, find the explicit wealth function $W(t)$.

\textbf{Problem 23 (Price Elasticity):} A commodity has a constant price elasticity of demand $E = -1.5$. Knowing that elasticity is given by $E = \frac{P}{Q}\frac{dQ}{dP}$, set up the differential equation and solve for the demand function $Q(P)$.

\textbf{Problem 24 (Vasicek Model):} An interest rate $r(t)$ reverts to a long-term mean of $0.05$ at a speed of $0.2$, modeled by $\frac{dr}{dt} = 0.2(0.05 - r)$. If the current interest rate is $r(0) = 0.01$, find the explicit solution for $r(t)$.

\textbf{Problem 25 (Population Modeling):} A bacterial population $P(t)$ grows according to the logistic equation $\frac{dP}{dt} = 0.05 P \left(1 - \frac{P}{500}\right)$. If the initial population is $P(0) = 100$, find the explicit equation for the population $P(t)$.

\subsection*{Part 7: Advanced Diagnostic Testing (''Find the Flaw'')}
For each problem below, a flawed solution is presented. Identify the exact mathematical error, explain why it is wrong in one sentence, and provide the correct calculation.

\textbf{Problem 26:} 
\textit{Question:} Solve $\frac{dy}{dx} = 3y$.
\textit{Flawed Solution:}
\begin{align*}
\frac{dy}{dx} &= 3y \\
\frac{1}{y} \, dy &= 3 \, dx \\
\ln|y| &= 3x \\
y &= e^{3x} + C 
\end{align*}

\textbf{Problem 27:} 
\textit{Question:} Solve $\frac{dy}{dx} = \frac{x}{y^2}$.
\textit{Flawed Solution:}
\begin{align*}
\frac{dy}{dx} &= \frac{x}{y^2} \\
y^{-2} \, dy &= x \, dx \\
\int y^{-2} \, dy &= \int x \, dx \\
\ln(y^2) &= \frac{1}{2}x^2 + C
\end{align*}

\textbf{Problem 28:} 
\textit{Question:} Determine if $\frac{dy}{dx} = 2x - xy$ is separable and separate it.
\textit{Flawed Solution:}
\begin{align*}
\frac{dy}{dx} &= 2x - xy \\
dy &= (2x - xy) dx \\
dy + xy \, dx &= 2x \, dx \\
\int dy + \int xy \, dx &= \int 2x \, dx
\end{align*}

\textbf{Problem 29:} 
\textit{Question:} Solve $\frac{dy}{dx} = x^2 y^3$.
\textit{Flawed Solution:}
\begin{align*}
\frac{dy}{dx} &= x^2 y^3 \\
\int y^3 \, dy &= \int x^2 \, dx \\
\frac{1}{4}y^4 &= \frac{1}{3}x^3 + C 
\end{align*}

\textbf{Problem 30:} 
\textit{Question:} Solve the IVP $\frac{dy}{dx} = -\frac{x}{y}, \quad y(0) = -5$.
\textit{Flawed Solution:}
\begin{align*}
y \, dy &= -x \, dx \\
\frac{1}{2}y^2 &= -\frac{1}{2}x^2 + C \\
y^2 &= -x^2 + C_1 \\
(-5)^2 &= -(0)^2 + C_1 \implies 25 = C_1 \\
y^2 &= 25 - x^2 \\
y &= \sqrt{25 - x^2}
\end{align*}

\newpage
% ==========================================
% DETAILED SOLUTIONS
% ==========================================
\section{Detailed Solutions}

\subsection*{Part 1: Identification and Setup Solutions}

\textbf{Solution 1:}
First, isolate the derivative:
\[ \frac{dy}{dx} = -4\ln(x-y) \]
To be separable, the right side must be expressible as $g(x)p(y)$. Because $x$ and $y$ are trapped inside the natural logarithm as a difference, logarithmic identities cannot uncouple them into a strict product. 
\textbf{Final Answer: No}

\textbf{Solution 2:}
We use the property of exponents, $e^{a-b} = e^a e^{-b}$, to rewrite the numerator:
\[ \frac{dy}{dx} = \frac{x^2 e^{3x} e^{-4y}}{x^4 + 1} \]
Now, group the $x$ terms and $y$ terms:
\[ \frac{dy}{dx} = \left(\frac{x^2 e^{3x}}{x^4 + 1}\right) \cdot (e^{-4y}) \]
This strictly matches the form $g(x)p(y)$.
\textbf{Final Answer: Yes, where $g(x) = \frac{x^2 e^{3x}}{x^4 + 1}$ and $p(y) = e^{-4y}$.}

\textbf{Solution 3:}
Rearrange the equation to isolate $\frac{dy}{dt}$:
\[ (ty^3 + 5y^3)dy = 7t \, dt \]
Factor out the common term $y^3$ on the left side:
\[ y^3(t + 5)dy = 7t \, dt \]
Divide by $dt$ and by $y^3(t+5)$:
\[ \frac{dy}{dt} = \left(\frac{7t}{t+5}\right) \cdot \left(\frac{1}{y^3}\right) \]
\textbf{Final Answer: Yes, where $g(t) = \frac{7t}{t+5}$ and $p(y) = \frac{1}{y^3}$.}

\subsection*{Part 2: General Implicit \& Explicit Solutions}

\textbf{Solution 4:}
Separate the variables by dividing by $x$ and multiplying by $dt$:
\[ \frac{1}{x} \, dx = 5t^4 \, dt \]
Integrate both sides:
\[ \int \frac{1}{x} \, dx = \int 5t^4 \, dt \]
\[ \ln|x| = t^5 + C \]
\textbf{Final Answer: $\ln|x| = t^5 + C$}

\textbf{Solution 5:}
Expand the exponential in the denominator:
\[ \frac{dy}{dt} = \frac{3}{y e^t e^{3y}} \]
Separate the variables by multiplying by $y e^{3y}$ and $dt$:
\[ y e^{3y} \, dy = 3e^{-t} \, dt \]
Integrate both sides. For the left side, use Integration by Parts. Let $u = y \implies du = dy$ and $dv = e^{3y} dy \implies v = \frac{1}{3}e^{3y}$:
\[ \int y e^{3y} \, dy = \frac{1}{3}y e^{3y} - \int \frac{1}{3}e^{3y} \, dy = \frac{1}{3}y e^{3y} - \frac{1}{9}e^{3y} \]
For the right side:
\[ \int 3e^{-t} \, dt = -3e^{-t} + C \]
Equate the integrals:
\textbf{Final Answer: $\frac{1}{3}y e^{3y} - \frac{1}{9}e^{3y} = -3e^{-t} + C$}

\textbf{Solution 6:}
Factor $t^2$ from the $dt$ term:
\[ t^2(1 + 4y^3) \, dt = -e^{t^3} y^2 \, dy \]
Separate the variables by dividing by $(1+4y^3)$ and $e^{t^3}$:
\[ \frac{y^2}{1+4y^3} \, dy = -t^2 e^{-t^3} \, dt \]
Integrate the left side using $u = 1+4y^3 \implies du = 12y^2 dy \implies \frac{1}{12}du = y^2 dy$:
\[ \int \frac{1}{12} \frac{1}{u} \, du = \frac{1}{12}\ln|1+4y^3| \]
Integrate the right side using $w = -t^3 \implies dw = -3t^2 dt \implies \frac{1}{3}dw = -t^2 dt$:
\[ \int \frac{1}{3} e^w \, dw = \frac{1}{3}e^{-t^3} + C \]
\textbf{Final Answer: $\frac{1}{12}\ln|1+4y^3| = \frac{1}{3}e^{-t^3} + C$}

\textbf{Solution 7:}
Separate variables by dividing by $\cos^2(y)$ and multiplying by $dx$:
\[ \frac{1}{\cos^2(y)} \, dy = 9x \, dx \]
Use the trigonometric identity $1/\cos^2(y) = \sec^2(y)$:
\[ \int \sec^2(y) \, dy = \int 9x \, dx \]
\[ \tan(y) = \frac{9}{2}x^2 + C \]
Take the inverse tangent to solve explicitly:
\textbf{Final Answer: $y = \arctan\left(\frac{9}{2}x^2 + C\right)$}

\textbf{Solution 8:}
Separate variables by dividing by $(y-2)$:
\[ \frac{1}{y-2} \, dy = x^4 \, dx \]
Integrate both sides:
\[ \ln|y-2| = \frac{1}{5}x^5 + C \]
Exponentiate both sides (Absolute value trick):
\[ |y-2| = e^{x^5/5 + C} = e^C \cdot e^{x^5/5} \]
Let $A = \pm e^C$:
\[ y - 2 = A e^{x^5/5} \]
\textbf{Final Answer: $y = 2 + A e^{x^5/5}$}

\textbf{Solution 9:}
Rearrange the equation:
\[ (3+x) \, dy = -\sqrt{y} \, dx \]
Separate variables:
\[ y^{-1/2} \, dy = -\frac{1}{3+x} \, dx \]
Integrate both sides:
\[ \int y^{-1/2} \, dy = \int -\frac{1}{3+x} \, dx \]
\[ 2\sqrt{y} = -\ln|3+x| + C \]
\textbf{Final Answer: $2\sqrt{y} + \ln|3+x| = C$}

\textbf{Solution 10:}
Factor $y$ from the numerator:
\[ \frac{dy}{dx} = \frac{y(x^2 - 2)}{x+5} \]
Separate variables:
\[ \frac{1}{y} \, dy = \frac{x^2 - 2}{x+5} \, dx \]
Use polynomial long division on the right side. $(x^2 - 2) \div (x+5) = x - 5$ with a remainder of $23$. Thus, $\frac{x^2-2}{x+5} = (x - 5) + \frac{23}{x+5}$.
Integrate both sides:
\[ \int \frac{1}{y} \, dy = \int \left( x - 5 + \frac{23}{x+5} \right) dx \]
\[ \ln|y| = \frac{1}{2}x^2 - 5x + 23\ln|x+5| + C \]
Exponentiate to isolate $y$, defining $A = \pm e^C$:
\[ y = A e^{\frac{1}{2}x^2 - 5x + 23\ln(x+5)} = A e^{\frac{1}{2}x^2 - 5x} \cdot e^{\ln((x+5)^{23})} \]
\textbf{Final Answer: $y = A (x+5)^{23} e^{\frac{1}{2}x^2 - 5x}$}

\subsection*{Part 3: Initial Value Problems (IVPs)}

\textbf{Solution 11:}
Separate variables by multiplying by $t \, dt$ and $\frac{y^3+1}{y}$:
\[ \frac{y^3+1}{y} \, dy = t \cos(t) \, dt \]
Split the fraction on the left: $(y^2 + \frac{1}{y}) \, dy = t \cos(t) \, dt$.
Integrate the left side: $\frac{1}{3}y^3 + \ln|y|$.
Integrate the right side by Parts ($u = t, dv = \cos(t) dt \implies du = dt, v = \sin(t)$):
\[ \int t \cos(t) \, dt = t\sin(t) - \int \sin(t) \, dt = t\sin(t) + \cos(t) + C \]
Combine:
\[ \frac{1}{3}y^3 + \ln|y| = t\sin(t) + \cos(t) + C \]
Apply initial condition $y(\pi) = 1$:
\[ \frac{1}{3}(1)^3 + \ln(1) = \pi\sin(\pi) + \cos(\pi) + C \]
\[ \frac{1}{3} + 0 = 0 - 1 + C \implies C = \frac{4}{3} \]
\textbf{Final Answer: $\frac{1}{3}y^3 + \ln|y| = t\sin(t) + \cos(t) + \frac{4}{3}$}

\textbf{Solution 12:}
Separate variables:
\[ \frac{1}{\cos^2(y)} \, dy = 2x \, dx \implies \sec^2(y) \, dy = 2x \, dx \]
Integrate both sides:
\[ \tan(y) = x^2 + C \]
Apply initial condition $y(1) = \frac{3\pi}{4}$:
\[ \tan\left(\frac{3\pi}{4}\right) = (1)^2 + C \implies -1 = 1 + C \implies C = -2 \]
Substitute $C$ back and solve for $y$:
\[ \tan(y) = x^2 - 2 \implies y = \arctan(x^2 - 2) + k\pi \]
Since the range of standard arctan is $(-\pi/2, \pi/2)$ and we need $y(1) = 3\pi/4$, we add $\pi$.
\textbf{Final Answer: $y = \arctan(x^2 - 2) + \pi$}

\textbf{Solution 13:}
Separate variables:
\[ \frac{1}{y+5} \, dy = x^3 \, dx \]
Integrate both sides:
\[ \ln|y+5| = \frac{1}{4}x^4 + C \]
Exponentiate:
\[ y + 5 = A e^{x^4/4} \implies y = -5 + A e^{x^4/4} \]
Apply initial condition $y(0) = 2$:
\[ 2 = -5 + A e^0 \implies A = 7 \]
\textbf{Final Answer: $y = -5 + 7e^{x^4/4}$}

\textbf{Solution 14:}
Separate variables:
\[ y \, dy = -x \, dx \]
Integrate both sides:
\[ \frac{1}{2}y^2 = -\frac{1}{2}x^2 + C_1 \implies y^2 + x^2 = C \]
Apply initial condition $x = 4, y = -3$:
\[ (-3)^2 + (4)^2 = C \implies 9 + 16 = 25 \implies C = 25 \]
Solve for $y$: $y^2 = 25 - x^2 \implies y = \pm\sqrt{25-x^2}$.
Since $y(4) = -3$ is negative, we MUST choose the negative branch.
\textbf{Final Answer: $y = -\sqrt{25-x^2}$}

\textbf{Solution 15:}
Separate variables:
\[ y^{-2} \, dy = t \ln(t) \, dt \]
Integrate the left side: $-y^{-1}$.
Integrate the right side by Parts ($u = \ln(t), dv = t \, dt \implies du = \frac{1}{t} dt, v = \frac{1}{2}t^2$):
\[ \int t \ln(t) \, dt = \frac{1}{2}t^2 \ln(t) - \int \frac{1}{2}t \, dt = \frac{1}{2}t^2 \ln(t) - \frac{1}{4}t^2 + C \]
Combine:
\[ -\frac{1}{y} = \frac{1}{2}t^2 \ln(t) - \frac{1}{4}t^2 + C \]
Apply initial condition $y(1) = -2$:
\[ -\frac{1}{-2} = \frac{1}{2}(1)^2 \ln(1) - \frac{1}{4}(1)^2 + C \]
\[ \frac{1}{2} = 0 - \frac{1}{4} + C \implies C = \frac{3}{4} \]
\textbf{Final Answer: $-\frac{1}{y} = \frac{1}{2}t^2 \ln(t) - \frac{1}{4}t^2 + \frac{3}{4}$}

\subsection*{Part 4: Singular Solutions}

\textbf{Solution 16:}
Separate variables by dividing by $y^{1/3}$:
\[ y^{-1/3} \, dy = dx \]
Integrate both sides:
\[ \frac{3}{2}y^{2/3} = x + C \implies y^{2/3} = \frac{2}{3}(x+C) \implies y = \pm\left(\frac{2}{3}(x+C)\right)^{3/2} \]
When we divided by $y^{1/3}$, we assumed $y \neq 0$. If we test $y=0$ in the original ODE, $\frac{dy}{dx} = 0$, and $(0)^{1/3} = 0$. So $y=0$ is a valid solution not covered by the general form.
\textbf{Final Answer: General solution $y = \pm\left(\frac{2}{3}(x+C)\right)^{3/2}$, Singular solution $y = 0$.}

\textbf{Solution 17:}
Separate variables:
\[ (y-5)^{-3/4} \, dy = x \, dx \]
Integrate:
\[ 4(y-5)^{1/4} = \frac{1}{2}x^2 + C \implies (y-5)^{1/4} = \frac{1}{8}x^2 + C_1 \implies y = 5 + \left(\frac{1}{8}x^2 + C_1\right)^4 \]
Dividing by $(y-5)^{3/4}$ invalidates $y=5$. Checking $y=5$ in the original ODE yields $0=0$.
\textbf{Final Answer: General solution $y = 5 + \left(\frac{1}{8}x^2 + C_1\right)^4$, Singular solution $y = 5$.}

\textbf{Solution 18:}
Separate variables:
\[ \frac{1}{16 - y^2} \, dy = dx \]
Use partial fractions: $\frac{1}{16-y^2} = \frac{1}{(4-y)(4+y)} = \frac{1/8}{4-y} + \frac{1/8}{4+y}$.
\[ \frac{1}{8} \int \left( \frac{1}{4+y} + \frac{1}{4-y} \right) dy = \int dx \]
\[ \frac{1}{8} \left( \ln|4+y| - \ln|4-y| \right) = x + C \implies \ln\left|\frac{4+y}{4-y}\right| = 8x + C_1 \]
Dividing by $16-y^2$ invalidates $y = \pm 4$. Both satisfy the original equation.
\textbf{Final Answer: General solution $\ln\left|\frac{4+y}{4-y}\right| = 8x + C_1$, Singular solutions $y = 4, y = -4$.}

\subsection*{Part 5: Logistic Growth \& Partial Fractions}

\textbf{Solution 19:}
Separate variables:
\[ \frac{1}{y(1 - y/10)} \, dy = 3 \, dt \implies \frac{10}{y(10 - y)} \, dy = 3 \, dt \]
Partial fractions decomposition:
\[ \frac{10}{y(10 - y)} = \frac{A}{y} + \frac{B}{10 - y} \implies 10 = A(10-y) + By \implies A=1, B=1 \]
Integrate:
\[ \int \left( \frac{1}{y} + \frac{1}{10 - y} \right) dy = \int 3 \, dt \]
\[ \ln|y| - \ln|10-y| = 3t + C \implies \ln\left|\frac{y}{10-y}\right| = 3t + C \]
\textbf{Final Answer: $\ln\left|\frac{y}{10-y}\right| = 3t + C$}

\textbf{Solution 20:}
Separate variables:
\[ \frac{1}{y^2 - 1} \, dy = \frac{1}{x} \, dx \]
Partial fractions for the left side:
\[ \frac{1}{y^2 - 1} = \frac{1}{(y-1)(y+1)} = \frac{1/2}{y-1} - \frac{1/2}{y+1} \]
Integrate:
\[ \frac{1}{2} \int \left( \frac{1}{y-1} - \frac{1}{y+1} \right) dy = \int \frac{1}{x} \, dx \]
\[ \frac{1}{2} (\ln|y-1| - \ln|y+1|) = \ln|x| + C \implies \ln\left|\frac{y-1}{y+1}\right| = 2\ln|x| + C_1 \]
\textbf{Final Answer: $\ln\left|\frac{y-1}{y+1}\right| = 2\ln|x| + C_1$}

\textbf{Solution 21:}
Factor the quadratic: $y^2 - 7y + 12 = (y-3)(y-4)$.
Separate variables:
\[ \frac{1}{(y-3)(y-4)} \, dy = dt \]
Partial fractions: $\frac{A}{y-3} + \frac{B}{y-4} \implies 1 = A(y-4) + B(y-3) \implies A=-1, B=1$.
\[ \int \left( \frac{1}{y-4} - \frac{1}{y-3} \right) dy = \int dt \]
\[ \ln|y-4| - \ln|y-3| = t + C \implies \ln\left|\frac{y-4}{y-3}\right| = t + C \]
Exponentiate:
\[ \frac{y-4}{y-3} = A e^t \implies y - 4 = A e^t (y - 3) \implies y - A e^t y = 4 - 3 A e^t \]
\textbf{Final Answer: $y = \frac{4 - 3Ae^t}{1 - Ae^t}$}

\subsection*{Part 6: Real-World Application and Modeling}

\textbf{Solution 22:}
The mathematical model is $\frac{dW}{dt} = 0.06W$.
Separate variables: $\frac{1}{W} dW = 0.06 \, dt$.
Integrate: $\ln|W| = 0.06t + C \implies W = A e^{0.06t}$.
Use initial condition $W(0) = 5000 \implies A = 5000$.
\textbf{Final Answer: $W(t) = 5000 e^{0.06t}$}

\textbf{Solution 23:}
Set up the model using $E = -1.5$:
\[ -1.5 = \frac{P}{Q}\frac{dQ}{dP} \]
Separate variables:
\[ \frac{1}{Q} \, dQ = -\frac{1.5}{P} \, dP \]
Integrate:
\[ \ln|Q| = -1.5\ln|P| + C \]
Exponentiate ($A = e^C$):
\[ Q(P) = e^{-1.5\ln|P| + C} = A e^{\ln(P^{-1.5})} = A P^{-1.5} \]
\textbf{Final Answer: $Q(P) = A P^{-1.5}$}

\textbf{Solution 24:}
Separate variables:
\[ \frac{1}{0.05 - r} \, dr = 0.2 \, dt \]
Integrate (use u-sub, $u = 0.05-r \implies du = -dr$):
\[ -\ln|0.05 - r| = 0.2t + C \implies \ln|0.05 - r| = -0.2t - C \]
Exponentiate:
\[ 0.05 - r = A e^{-0.2t} \implies r(t) = 0.05 - A e^{-0.2t} \]
Apply initial condition $r(0) = 0.01$:
\[ 0.01 = 0.05 - A \implies A = 0.04 \]
\textbf{Final Answer: $r(t) = 0.05 - 0.04e^{-0.2t}$}

\textbf{Solution 25:}
Separate variables:
\[ \frac{1}{P(1 - P/500)} \, dP = 0.05 \, dt \implies \frac{500}{P(500-P)} \, dP = 0.05 \, dt \]
Partial fractions gives $\frac{1}{P} + \frac{1}{500-P}$.
\[ \int \left(\frac{1}{P} + \frac{1}{500-P}\right) dP = \int 0.05 \, dt \implies \ln|P| - \ln|500-P| = 0.05t + C \]
\[ \ln\left|\frac{P}{500-P}\right| = 0.05t + C \implies \frac{P}{500-P} = A e^{0.05t} \]
Apply initial condition $P(0) = 100$:
\[ \frac{100}{400} = A e^0 \implies A = \frac{1}{4} \]
Solve for $P$:
\[ P = \frac{1}{4}e^{0.05t}(500 - P) \implies P \left(1 + \frac{1}{4}e^{0.05t}\right) = 125e^{0.05t} \implies P = \frac{125e^{0.05t}}{1 + 0.25e^{0.05t}} = \frac{500}{4e^{-0.05t} + 1} \]
\textbf{Final Answer: $P(t) = \frac{500}{1 + 4e^{-0.05t}}$}

\subsection*{Part 7: Advanced Diagnostic Testing (''Find the Flaw'')}

\textbf{Solution 26:}
\textbf{The Flaw:} The constant of integration $C$ was added outside as $+C$ after exponentiating both sides, rather than absorbing it as a coefficient immediately after integration.
\textbf{Correct Step:} $\ln|y| = 3x + C \implies |y| = e^{3x+C} \implies y = A e^{3x}$.

\textbf{Solution 27:}
\textbf{The Flaw:} The student incorrectly integrated $y^{-2}$ using a natural logarithm instead of the power rule.
\textbf{Correct Step:} $\int y^{-2} \, dy = \frac{y^{-1}}{-1} = -\frac{1}{y}$. Thus, $-\frac{1}{y} = \frac{1}{2}x^2 + C$.

\textbf{Solution 28:}
\textbf{The Flaw:} The student attempted to separate variables using addition, which is algebraically invalid since integrals cannot separate variables over sums like $dy + xy \, dx$.
\textbf{Correct Step:} Factor first: $\frac{dy}{dx} = x(2 - y) \implies \frac{1}{2-y} \, dy = x \, dx$.

\textbf{Solution 29:}
\textbf{The Flaw:} The student multiplied by $y^3$ to move it to the left side, but algebraically they needed to divide by $y^3$.
\textbf{Correct Step:} $\frac{1}{y^3} \, dy = x^2 \, dx \implies \int y^{-3} \, dy = \int x^2 \, dx \implies -\frac{1}{2}y^{-2} = \frac{1}{3}x^3 + C$.

\textbf{Solution 30:}
\textbf{The Flaw:} When taking the square root in the final step, the student wrongly assumed the positive root branch, contradicting the negative initial condition $y(0) = -5$.
\textbf{Correct Step:} $y = \pm\sqrt{25 - x^2}$. Since $y(0)$ must be negative, the correct branch is $y = -\sqrt{25 - x^2}$.

\newpage
% ==========================================
% CONCEPT CHECKLIST RECAP
% ==========================================
\section{Concept Checklist \& Question Mapping (Recap)}
As requested, here is the mapping so you can easily reference problem numbers by their target concepts:

\begin{itemize}
    \item \textbf{Identification/Setup:} Q1, Q2, Q3
    \item \textbf{General Solutions:} Q4, Q5, Q6, Q7, Q8, Q9, Q10
    \item \textbf{Initial Value Problems:} Q11, Q12, Q13, Q14, Q15
    \item \textbf{Singular Solutions:} Q16, Q17, Q18
    \item \textbf{Partial Fractions / Logistic Growth:} Q19, Q20, Q21, Q25
    \item \textbf{Real-World Modeling:} Q22, Q23, Q24, Q25
    \item \textbf{Find the Flaw (Diagnostics):} Q26, Q27, Q28, Q29, Q30
\end{itemize}

\textbf{Algebraic Manipulations Addressed:}
\begin{itemize}
    \item Splitting Exponents: Q2, Q5
    \item Factoring: Q3, Q6, Q8, Q10
    \item Integration by Parts: Q5, Q11, Q15
    \item Trig Identities: Q7, Q12
    \item Logarithmic Exponentiation / Absolute Value Trick: Q8, Q10, Q13, Q19, Q24, Q25
    \item Negative Square Root Branches: Q14, Q30
\end{itemize}

\end{document}